\section{Fit With the Presidential Hackathon}\label{sec:fit}

The Presidential Hackathon is not a product competition; it is the Taiwanese
government's \textbf{policy-adoption pipeline}. Winning teams present to the
President; five Teams of Excellence per year receive government guidance,
public-private matchmaking, field-validation support, and mandatory follow-up
benefit tracking; past winners returned to the 2026 launch event to present
post-victory implementation progress, and at least one past project matured
into national policy\cite{pmc,moda}. The Handbook's own application form
contains ``development stage,'' ``existing validation,'' and ``future plans''
fields designed for early-stage projects\cite{handbook}. \textbf{Full
deployment is therefore not expected at entry --- a working proof of concept
plus a government-connected validation path is precisely the intended shape},
and claiming more would score worse under Feasibility (40\%), the heaviest
preliminary criterion.

\begin{table}[htbp]
\caption{Mapping to the official evaluation criteria\cite{handbook}}
\label{tab:criteria}
\small
\begin{tabularx}{\textwidth}{@{}p{0.28\textwidth}X@{}}
\toprule
\textbf{Criterion (prelim / final)} & \textbf{How CarbonPass answers it} \\
\midrule
Feasibility (40\% / 30\%) & Built entirely on open, named datasets (Table~\ref{tab:datasets}); no hardware, no partnerships required for the PoC; staged plan matched to the official calendar (Table~\ref{tab:plan}) \\
Innovation (30\% / 20\%) & First tool converting SME-native documents into EU-format product-level emissions data; carbon-aware scheduling driven by Taiwan's open 10-minute grid feed; capability audit (Table~\ref{tab:gap}) shows no existing overlap \\
Social Impact (30\% / 20\%) & Protects a nameable regional blue-collar community (2,600 firms; a top-3 world export industry) from a dated, quantified threat; inclusion delivered in the interface itself \\
Implementation \& Verification (--- / 30\%) & Pilot factory during mentorship; measured before/after ledger; weekly documented progress between submission and finals \\
\bottomrule
\end{tabularx}
\end{table}

On the theme: the Handbook defines digital inclusion as enabling \emph{``people
of different ages, regions, and needs to access technology and participate in
digital services on equal footing''}\cite{handbook}. CarbonPass maps to all
three axes without stretching the text --- \textbf{ages} (aging, family-run
firms asked to operate EU-grade data infrastructure overnight),
\textbf{regions} (one nameable southern district whose economic identity is at
stake), and \textbf{needs} (participation in what is now a de facto digital
service: carbon-verified market access). Historical precedent supports
economically framed winners: procurement transparency (2020), agricultural
carbon reduction (2022), a blockchain carbon wallet (2024), agri-AI
(2025)\cite{moda}.

\section{Honest Limits and Risks}\label{sec:risks}

\begin{itemize}[leftmargin=1.3em]
\item \textbf{The tool prepares verification; it does not replace it.} At
least 80\% of reported emissions must be verified by an accredited third
party\cite{omm}; CarbonPass's output is the structured input to that process,
never a certificate.
\item \textbf{The 2028 extension is pending legislation} until the European
Parliament and Council conclude; the proposal books no benefit from it, only a
tracked roadmap event\cite{akin}.
\item \textbf{Preferential fee rates belong to large emitters.} The SME's
incentives are customer retention, CBAM cost pass-through, energy savings, and
readiness --- the proposal says exactly that, nothing more\cite{moenvfee}.
\item \textbf{Brussels risk cuts both ways.} Further CBAM simplification could
shrink the compelled population; the extension could triple it. The beachhead
(live compulsion today) is chosen precisely to be robust to either outcome.
\item \textbf{The addressable core is 2,600 firms today}\cite{tpt-cbam} --- a
deliberately honest number. The ``big idea'' is the expansion path; the demo
scope is one cluster. Under a rubric that weights Feasibility at 40\%, honesty
is the competitive strategy.
\end{itemize}
